\ifdefined\directlua
  \documentclass[a4paper,11pt]{ltjsarticle}
\else
  \documentclass[a4paper,11pt,dvipdfmx]{jsarticle}
\fi
\usepackage[margin=20mm]{geometry}
\usepackage{KKchemstruct}

\newcommand{\bs}{\char`\\}
\newcommand{\bl}{\char`\{}
\newcommand{\br}{\char`\}}
\newcommand{\demo}[2]{\par\medskip\noindent\texttt{#1}\par\smallskip\noindent#2\par}

\begin{document}

\section*{KKchemstruct regression test}

\subsection*{1. ベンゼン環（\texttt{\bs KKbenzene}）}

\demo{\bs KKbenzene[]}{\KKbenzene[]}
\demo{\bs KKbenzene[2=OH,3=COOH]}{\KKbenzene[2=OH,3=COOH]}
\demo{\bs KKbenzene[1=COOH,4=COOH]}{\KKterephthalicacid}
\demo{\bs KKbenzene[1=OH,kekule=b]}{\KKbenzene[1=OH,kekule=b]}
\demo{\bs KKbenzene[2=OH,angle=30,sep=1.0em]}{\KKbenzene[2=OH,angle=30,sep=1.0em]}

\subsection*{2. 本文中に置く横向きの環}

本文の中に \KKphenylring{} を置いても行送りが崩れないこと、
\KKphenyl{COOH} のように右へ置換基が出ること、
鎖の途中に \ck{CH_3(CH_2)_{11}}\KKphenylene{SO_3H} と挟めることを確認します。

\subsection*{3. 示性式（\texttt{\bs KKchemchain}）}

\demo{\bs KKchemchain\bl R\_1-C\bs KKcarbonyl-OH\br}{\KKestergeneric}
\demo{\bs KKglycerol}{\KKglycerol}
\demo{\bs KKtriglyceride\bl R\_1\br\bl R\_2\br\bl R\_3\br}{\KKtriglyceride{R_1}{R_2}{R_3}}
\demo{\bs KKsulfuricacid（配位結合）}{\KKsulfuricacid}
\demo{\bs KKnitricacid}{\KKnitricacid}
\demo{\bs KKphosphoricacid}{\KKphosphoricacid}

\subsection*{4. 縮合環・幾何異性体}

\demo{\bs KKphthalicanhydride}{\KKphthalicanhydride}
\demo{\bs KKmaleicanhydride}{\KKmaleicanhydride}
\demo{\bs KKmaleicacid / \bs KKfumaricacid}%
  {\KKmaleicacid\hspace{4em}\KKfumaricacid}

\subsection*{5. 名称・注釈}

\demo{\bs KKname}{\KKname{\KKsalicylicacid}{サリチル酸}}

\demo{\bs KKchemmark（脱離部分の点線囲み）}{%
  \KKchemchain{R-C\KKcarbonyl-@{s}OH}\ +\ \KKchemchain{@{e}H-O-R'}%
  \KKchemmark{s}{e}}

\demo{\bs KKhbond（分子内水素結合）}{%
  \KKbenzene[2=@{oh}OH,3=C@{co}OOH]\KKhbond{oh}{co}}

\demo{\bs KKchemcross}{\KKchemcross{\KKphenol}}

\subsection*{6. 反応式}

\KKscheme{%
  \KKname{\KKsalicylicacid}{サリチル酸}
  \arrow{->[\ck{CH_3OH}][エステル化]}
  \KKname{\KKmethylsalicylate}{サリチル酸メチル}}

\KKbranchscheme
  {\KKname{\KKsalicylicacid}{サリチル酸}}
  {\ck{CH_3OH}}{エステル化}{\KKname{\KKmethylsalicylate}{サリチル酸メチル}}
  {\ck{(CH_3CO)_2O}}{アセチル化}{\KKname{\KKaspirin}{アセチルサリチル酸}}

\subsection*{7. 寸法の一括変更（\texttt{\bs KKchemsetup}）}

\begingroup
\KKchemsetup{ring sep=2.0em,bond width=0.09em,double bond sep=0.4em}
\KKsalicylicacid\hspace{4em}\KKphenyl{COOH}
\endgroup

変更はグループ内に閉じているので、ここでは元の寸法に戻ります: \KKsalicylicacid

\subsection*{8. 化学式マクロ}

\ck{CH_3COOH}、\ck{RCOO\ckm}、\ck{Na\ckp} と、
数式モード中の $\ck{H_2O}$ が同じ体裁で組まれること。

\end{document}
